\documentclass[11pt]{article}
\usepackage[margin=1.1in]{geometry}
\usepackage{amsmath,amssymb,amsthm,mathtools,booktabs,array,tabularx,longtable}
\usepackage[dvipsnames]{xcolor}
\usepackage[colorlinks=true,linkcolor=MidnightBlue,citecolor=MidnightBlue,urlcolor=MidnightBlue]{hyperref}
\usepackage{fancyhdr}
\pagestyle{fancy}\fancyhf{}\renewcommand{\headrulewidth}{0pt}
\fancyfoot[R]{\thepage}
\fancyfoot[L]{\scriptsize\color{gray}Generated by the Transport Lab from approved parts of its store; not a deposited paper. Papers of record: \href{https://zenodo.org/communities/transport-papers}{Zenodo, Transport Papers}}
\fancypagestyle{plain}{\fancyhf{}\fancyfoot[R]{\thepage}\fancyfoot[L]{\scriptsize\color{gray}Generated by the Transport Lab from approved parts of its store; not a deposited paper. Papers of record: \href{https://zenodo.org/communities/transport-papers}{Zenodo, Transport Papers}}\renewcommand{\headrulewidth}{0pt}}
\newtheorem{theorem}{Theorem}[section]
\newtheorem{proposition}[theorem]{Proposition}
\newtheorem{lemma}[theorem]{Lemma}
\newtheorem{corollary}[theorem]{Corollary}
\newtheorem{conjecture}[theorem]{Conjecture}
\newtheorem{refutation}[theorem]{Refutation}
\newtheorem{prediction}[theorem]{Prediction}
\newtheorem{principle}[theorem]{Principle}
\newtheorem{observation}[theorem]{Observation}
\newtheorem{axiom}[theorem]{Axiom}
\theoremstyle{definition}
\newtheorem{definition}[theorem]{Definition}
\newtheorem{fact}[theorem]{Fact}
\newtheorem{assembled}[theorem]{Assembly}
\newtheorem{postulate}[theorem]{Postulate}
\newtheorem{premise}[theorem]{Premise}
\newtheorem{hypothesis}[theorem]{Hypothesis}
\newtheorem{remark}[theorem]{Remark}
\newtheorem{example}[theorem]{Example}
\theoremstyle{plain}
\setlength{\marginparwidth}{0.8in}
\newcommand{\changedmark}[2]{\marginpar{\raggedright\scriptsize\color{#1}#2}}
\providecommand{\Aa}{}\renewcommand{\Aa}{\mathcal{A}}
\providecommand{\Ad}{}\renewcommand{\Ad}{\widetilde{\mathrm{Ad}}}
\providecommand{\Aev}{}\renewcommand{\Aev}{\mathcal{A}^{+}}
\providecommand{\Aod}{}\renewcommand{\Aod}{\mathcal{A}^{-}}
\providecommand{\B}{}\renewcommand{\B}{\mathrm{B}}
\providecommand{\C}{}\renewcommand{\C}{\mathbb{C}}
\providecommand{\CAR}{}\renewcommand{\CAR}{\mathrm{CAR}}
\providecommand{\CS}{}\renewcommand{\CS}{{\mathbb{S}^{+}}}
\providecommand{\CSm}{}\renewcommand{\CSm}{{\mathbb{S}^{-}}}
\providecommand{\Cl}{}\renewcommand{\Cl}{\mathrm{Cl}}
\providecommand{\Con}{}\renewcommand{\Con}{\mathcal{C}}
\providecommand{\Dl}{}\renewcommand{\Dl}{\Delta}
\providecommand{\EN}{}\renewcommand{\EN}{E_{N}}
\providecommand{\End}{}\renewcommand{\End}{\mathrm{End}}
\providecommand{\F}{}\renewcommand{\F}{\mathbb{F}}
\providecommand{\G}{}\renewcommand{\G}{\Gamma}
\providecommand{\HH}{}\renewcommand{\HH}{\mathbb{H}}
\providecommand{\IS}{}\renewcommand{\IS}{I\sigma_3}
\providecommand{\J}{}\renewcommand{\J}{\mathbf{J}}
\providecommand{\LH}{}\renewcommand{\LH}{L_{H}}
\providecommand{\Mat}{}\renewcommand{\Mat}{\mathrm{M}}
\providecommand{\Mo}{}\renewcommand{\Mo}{M^{\circ}}
\providecommand{\Mp}{}\renewcommand{\Mp}{M^{+}}
\providecommand{\Msun}{}\renewcommand{\Msun}{M_{\odot}}
\providecommand{\Mt}{}\renewcommand{\Mt}{\widetilde{M}}
\providecommand{\Ncone}{}\renewcommand{\Ncone}{\mathcal{N}}
\providecommand{\Neff}{}\renewcommand{\Neff}{N_{\mathrm{eff}}}
\providecommand{\PCS}{}\renewcommand{\PCS}{\mathbb{P}{\mathbb{S}^{+}}}
\providecommand{\PT}{}\renewcommand{\PT}{\mathbb{PT}}
\providecommand{\Pin}{}\renewcommand{\Pin}{\mathrm{Pin}}
\providecommand{\Q}{}\renewcommand{\Q}{\mathbb{Q}}
\providecommand{\R}{}\renewcommand{\R}{\mathbb{R}}
\providecommand{\Ric}{}\renewcommand{\Ric}{\mathrm{Ric}}
\providecommand{\SBH}{}\renewcommand{\SBH}{S_{\mathrm{BH}}}
\providecommand{\Sc}{}\renewcommand{\Sc}{\mathbb{S}}
\providecommand{\Sch}{}\renewcommand{\Sch}{\mathrm{P}}
\providecommand{\Sig}{}\renewcommand{\Sig}{\Sigma}
\providecommand{\Som}{}\renewcommand{\Som}{S_\omega}
\providecommand{\Sp}{}\renewcommand{\Sp}{\mathbb{S}}
\providecommand{\Sph}{}\renewcommand{\Sph}{\mathrm{S}}
\providecommand{\Spi}{}\renewcommand{\Spi}{S_\pi}
\providecommand{\Spin}{}\renewcommand{\Spin}{\mathrm{Spin}}
\providecommand{\Sth}{}\renewcommand{\Sth}{S_{\mathrm{th}}}
\providecommand{\T}{}\renewcommand{\T}{\mathbb{T}}
\providecommand{\TH}{}\renewcommand{\TH}{T_{\mathrm{H}}}
\providecommand{\Tr}{}\renewcommand{\Tr}{\operatorname{Tr}}
\providecommand{\UHF}{}\renewcommand{\UHF}{\mathrm{UHF}}
\providecommand{\Z}{}\renewcommand{\Z}{\mathbb{Z}}
\providecommand{\Zc}{}\renewcommand{\Zc}{\mathrm{Z}}
\providecommand{\cmg}{}\renewcommand{\cmg}{\mathrm{cm^{2}\,g^{-1}}}
\providecommand{\dS}{}\renewcommand{\dS}{\mathrm{dS}}
\providecommand{\dsc}{}\renewcommand{\dsc}{\mathrm{dsc}}
\providecommand{\dtPG}{}\renewcommand{\dtPG}{dt}
\providecommand{\dtS}{}\renewcommand{\dtS}{dt_{\mathrm{S}}}
\providecommand{\ellv}{}\renewcommand{\ellv}{\ell}
\providecommand{\er}{}\renewcommand{\er}{e_r}
\providecommand{\et}{}\renewcommand{\et}{e_t}
\providecommand{\ev}{}\renewcommand{\ev}{\mathrm{ev}}
\providecommand{\gl}{}\renewcommand{\gl}{\mathfrak{gl}}
\providecommand{\hatV}{}\renewcommand{\hatV}{\widehat{V}}
\providecommand{\hb}{}\renewcommand{\hb}{\bar{h}}
\providecommand{\hf}{}\renewcommand{\hf}{\hat{f}}
\providecommand{\id}{}\renewcommand{\id}{\mathrm{id}}
\providecommand{\kpc}{}\renewcommand{\kpc}{\mathrm{kpc}}
\providecommand{\lP}{}\renewcommand{\lP}{\ell_{P}}
\providecommand{\mD}{}\renewcommand{\mD}{\mathbf{m}_{D}}
\providecommand{\mb}{}\renewcommand{\mb}{\bar m}
\providecommand{\mnu}{}\renewcommand{\mnu}{\Sigma m_{\nu}}
\providecommand{\nInf}{}\renewcommand{\nInf}{n_\infty}
\providecommand{\nO}{}\renewcommand{\nO}{n_0}
\providecommand{\ninf}{}\renewcommand{\ninf}{n_{\infty}}
\providecommand{\nz}{}\renewcommand{\nz}{n_{0}}
\providecommand{\proj}{}\renewcommand{\proj}{\operatorname{proj}}
\providecommand{\pv}{}\renewcommand{\pv}{\pi}
\providecommand{\rad}{}\renewcommand{\rad}{\mathrm{rad}}
\providecommand{\sgn}{}\renewcommand{\sgn}{\operatorname{sgn}}
\providecommand{\shA}{}\renewcommand{\shA}{\mathcal{A}}
\providecommand{\shC}{}\renewcommand{\shC}{\mathcal{C}}
\providecommand{\shE}{}\renewcommand{\shE}{\mathcal{E}}
\providecommand{\shJ}{}\renewcommand{\shJ}{\mathcal{J}}
\providecommand{\shK}{}\renewcommand{\shK}{\mathcal{K}}
\providecommand{\shL}{}\renewcommand{\shL}{L}
\providecommand{\shW}{}\renewcommand{\shW}{\mathcal{W}}
\providecommand{\sll}{}\renewcommand{\sll}{\mathfrak{sl}}
\providecommand{\so}{}\renewcommand{\so}{\mathfrak{so}}
\providecommand{\su}{}\renewcommand{\su}{\mathfrak{su}}
\providecommand{\supp}{}\renewcommand{\supp}{\operatorname{supp}}
\providecommand{\tr}{}\renewcommand{\tr}{\operatorname{tr}}
\providecommand{\uu}{}\renewcommand{\uu}{\mathfrak{u}}
\providecommand{\br}{}\renewcommand{\br}[1]{\langle #1\rangle}
\providecommand{\bra}{}\renewcommand{\bra}[1]{\langle #1|}
\providecommand{\file}{}\renewcommand{\file}[1]{\texttt{#1}}
\providecommand{\hp}{}\renewcommand{\hp}[2]{( #1 \,|\, #2 )}
\providecommand{\ip}{}\renewcommand{\ip}[2]{\langle #1, #2\rangle}
\providecommand{\ket}{}\renewcommand{\ket}[1]{|#1\rangle}
\providecommand{\rel}{}\renewcommand{\rel}[2]{D\bigl(#1 \,\|\, #2\bigr)}
\providecommand{\rev}{}\renewcommand{\rev}[1]{\widetilde{#1}}
\providecommand{\slashed}{}\renewcommand{\slashed}[1]{{#1}\mkern-10mu/}
\title{\bfseries Selecting the State: Parity, Regularity, and the Irreducible Calibration behind the Hawking Temperature: A companion note to Transport Papers T5, Step 4d}
\author{Leon Butler\\ \normalsize\href{mailto:lbutler.math@gmail.com}{\texttt{lbutler.math@gmail.com}}\\ \normalsize ORCID: \href{https://orcid.org/0009-0006-4449-5617}{0009-0006-4449-5617}}
\date{Generated by the Transport Lab, 7 November 2026}
\begin{document}
\maketitle
\begin{abstract}
\normalsize
Step 4d of T5 proves thermality at the Hawking temperature for any state that is pure, boost invariant and entangling, and declares its status plainly as ``up to a state''. This note asks what selects the state, and answers in four parts. First, the descent map of the chain is identified: it is exactly the fermion parity average, so the metric level is the parity even level, and for every state whatsoever the descended single mode marginal is diagonal, hence Gibbs at some inverse temperature. Per mode thermality is therefore not a property that can fail; the entire physical content of Step 4d is one number per mode. Second, under the univalence premise the ascended marginals are already diagonal, so the parity protection lemma of the finite ledger model is a theorem about all states and hypothesis (ii) is discharged. Third, the residual number is pinned by a rigidity theorem: for each inverse temperature there is exactly one KMS state, no flow preserving automorphism relates two of them, and the only algebraic relation available is the orientation flip. The temperature is therefore not algebra intrinsic, and no condition stated in the algebra and its flow can fix it. Fourth, one named regularity premise fixes it, and the framework supplies the premise's habitat: the curvature selected eigenplane pair of T5 Theorem 6.1 contains both the Lorentzian boost and its Euclidean partner, whose rotor has period $4\pi$ and whose conjugation action has period $2\pi$, giving $T = \kappa/2\pi$ and, as a by-product, the antiperiodicity responsible for Fermi-Dirac rather than Bose-Einstein statistics. A fifth result cuts the other way: exact KMS for the geometric flow forces a product state, hence no inter-quanta correlations at all, so the correlations carrying the information in Steps 5 and 6b are precisely the deviation from exact stationarity, and they cannot be selected by any equilibrium condition. Purity is shown to be the canonical GNS purification rather than an assumption. Step 4d's status changes from ``up to a state'' to ``proved from univalence and one regularity premise''; the residual state dependence migrates to Step 6b, where the chain already prices it.
\end{abstract}
\medskip
\noindent\textbf{MSC 2020.} Primary 83C57, 46L60; Secondary 81T05, 82B10.\\
\textbf{Keywords.} state selection; Hawking temperature; KMS states; modular theory; fermion parity superselection; Hartle-Hawking state; conditional expectation; Clifford algebra.
\newpage
\noindent \textbf{Authorship and disclosure. }The founding idea of the series is the author's: that division by zero in a Clifford algebra should be read not as a prohibition but as a transport into a higher-dimensional algebra, and that the transported division need not be accompanied by an inverse operation. The direction of investigation and the critical review of the results are likewise the author's. The mathematical development and the drafting of the text were carried out in extended collaboration with Claude Fable 5 (Anthropic), an AI system.

\noindent\fbox{\parbox{\dimexpr\linewidth-2\fboxsep-2\fboxrule}{\textbf{This is not a deposited paper.} It was generated by the Transport Lab on 7 November 2026 and assembled from approved parts of its store. It follows results across several papers of the series; each is stated with its proof and credited to the paper that first proved it. The official versions of those papers are on Zenodo, in the Transport Papers community, \url{https://zenodo.org/communities/transport-papers}.}}
\newpage

\section{The gap, stated exactly}

\phantomsection\label{c:note.state-selection.setting-gap-stated-exactly}
Step 4d of T5 \cite{ButlerT5} proves the following. Let the completed ladder pair $(v, v^{*})$ of the singular locus carry a state that is (i) pure, (ii) invariant under the boost flow generated by $\sigma$, and (iii) entangling across the two ends of the pair, that is cyclic and separating for the one-sided subalgebra. Then the descended state is KMS for the boost flow at modular temperature $2\pi$ of rapidity, which by the geometric calibration of \cite[Prop.~6.2]{ButlerT5} is
\[
T \;=\; \frac{\kappa}{2\pi} \;=\; \frac{1}{8\pi M} \;=\; \TH ,
\]
with per-mode reduced density matrix the Fermi-Dirac matrix. The theorem is stated with its converse: a state violating (ii) or (iii) reduces non-thermally under the same projector. The status line in the chain table reads \emph{up to a state}, and the paper says in as many words that the projection alone cannot select a temperature and that honesty forbids pretending otherwise.

That honesty is the point of departure here. Every derivation of the Hawking effect meets the same wall. Hawking's original calculation \cite{Hawking1975} selects the state by the collapse history and the choice of an ``in'' vacuum; Hartle and Hawking \cite{HartleHawking1976} select it by regularity of a Euclidean path integral; Unruh \cite{Unruh1976} by regularity on the future horizon; Boulware's alternative \cite{Boulware1975} is the same field theory with a different state and radiates nothing; Kay and Wald \cite{KayWald1991} prove, in the rigorous setting, that a stationary nonsingular quasi-free state on a spacetime with a bifurcate Killing horizon is unique when it exists, and is KMS at the Hawking temperature, which is exactly a theorem about the strength of the regularity premise; and Fredenhagen and Haag \cite{FredenhagenHaag1990} localise the input to the short-distance behaviour of the state near the horizon. The wall is not an accident of any one formalism, and it is not a wall this note claims to knock down.

What this note does claim is that the wall stands in a different place than Step 4d's three hypotheses suggest, and that once it is put in the right place, two of the three hypotheses are theorems, the third is an output, and what remains is one real number and one named premise. The results, with statuses, are in the table accompanying \ref{c:singularity.descent-parity-average}.

 \begin{tabular}{l>{ \arraybackslash}p{6.4cm}>{ \arraybackslash}p{4.6cm}l}
\toprule
Result & Content & Consequence for T5 & Status \\
\midrule
Thm \ref{c:singularity.descent-parity-average} & the descent $E$ is the fermion parity average; the metric level is the parity even level & explains \cite[Fact 5.9]{ButlerT5} structurally & proved \\
Thm \ref{c:singularity.descended-form-state-independent} & for \emph{every} state, the descended single-mode marginal is diagonal, hence Gibbs at some $\beta_k$ & the \emph{form} in Thm 7.3 needs no hypothesis & proved \\
Thm \ref{c:singularity.ascended-marginals-are-diagonal} & under univalence (P), ascended marginals are already diagonal & hypothesis (ii) discharged; \cite[Lem.~12.1]{ButlerT5} generalised & proved from (P) \\
Thm \ref{c:singularity.kms-uniqueness-factor} & exactly one KMS state per $\beta$; $\rho \propto e^{-\beta K}$ & the state class is a line, not a class & proved \\
Thm \ref{c:singularity.rate-rigidity-temperature-not-algebra-intrinsic} & no flow-preserving automorphism relates $\beta$ to $\beta'$; the only relation is $\beta \mapsto -\beta$ & the temperature is not algebra intrinsic & proved \\
Thm \ref{c:singularity.exact-equilibrium-over-determines-boundary-modes} & exact KMS for the geometric flow forces a product state, hence zero inter-quanta correlation & equilibrium cannot supply Step 6b's correlations & proved \\
Prop \ref{c:state-selection.euclidean-partner-curvature-selected} & the Euclidean partner of the boost is curvature selected; its rotor has period $4\pi$ & the habitat of premise (R) is internal & proved \\
Thm \ref{c:state-selection.rate-discharge-hypothesis-iii} & under (R), $\beta = 2\pi$ in rapidity, so $T = \kappa/2\pi$; $0 < \beta < \infty$ discharges (iii) & the calibration enters once & proved from (R) \\
Cor \ref{c:state-selection.where-fermi-dirac-comes} & the $2\pi$ rotor is $-1$: the odd sector is antiperiodic & the $+1$ in Fermi-Dirac, structurally & proved \\
Thm \ref{c:state-selection.partner-purification-are-canonical} & the two-sided pure state is the GNS purification of the one-sided state & hypothesis (i) is an output & proved \\
Prop \ref{c:singularity.number-diagonal-dynamics-models-scrambling-derived} & ledger locality forces number-diagonal internal dynamics & the model's diagonal scrambling derived & from (L$_E$) \\
\bottomrule
\end{tabular}

\emph{The results of this note and what each does to Step 4d's hypotheses. (P) is the univalence premise of the section ``The descent is the parity map''; (R) the regularity premise of the section ``Fixing the rate: one premise, and its habitat inside the framework''; (L$_E$) the ledger locality premise of the section ``What equilibrium cannot supply: correlations''.}

The one-sentence summary is this. Thermality of a single fermionic mode is not a property that a state can fail to have, so the state-selection problem in Step 4d is not about the shape of the reduced density matrix at all; it is about one number, the number is a ratio of two clocks, and a ratio of clocks requires a calibration that no amount of algebra can supply. The chain's contribution is that it says exactly which two clocks, and proves that the calibration enters exactly once.

\subsection*{Notation and standing conventions}
\phantomsection\label{c:note.state-selection.setting-notation-standing-conventions}
Throughout, $\Aa$ denotes the completed ladder algebra of the singular locus: at a single point of the blown-up singular boundary $E \cong \R \times S^{2}$ the algebra generated by the Rule 1 pair $(v, v^{*})$ with
\begin{equation}
v^{2} = (v^{*})^{2} = 0, \qquad v v^{*} + v^{*} v = 1 ,
\end{equation}
so that $\Aa \cong \Mat_2(\C)$ as a factor, and over all of $E$ the CAR algebra generated by one such pair per point \cite[Thm.~8.7, Thm.~8.12]{ButlerT4}. We write $N = v^{*}v$ for the number operator, $P = N$ for the occupation projector, $\sigma = v^{*}v - vv^{*} = 2N - 1$ for the modular bivector of \cite[Thm.~6.3]{ButlerT5}, and $K = \sigma/2 = N - \tfrac12$ for the boost generator. The descent is
\begin{equation}
E(x) \;=\; PxP + (1-P)\,x\,(1-P)
\end{equation}
of \cite[Prop.~5.5]{ButlerT5}, and $\EN$ its multi-mode composition \cite[Def.~9.1]{ButlerT5}. Mode labels $k$ run over the ledger; $\omega_k > 0$ is the mode's boost weight. All statements about $\Aa$ at a single point are finite dimensional and are proved outright; the statements that need the infinite system are flagged where they occur. Units and curvature conventions are T5's.

\section*{What we assume}
The results below are taken as given in this paper; each result that rests on one says so.
\begin{postulate}[(L$_E$): boundary-mode locality, as descent covariance on finite windows]\label{c:singularity.l-boundary-mode-locality}
\emph{Premise (L$_{E}$), boundary-mode locality.} On every finite window of the boundary modes, $v_{1}, \dots, v_{n}$ satisfying the CAR, the internal dynamics $\alpha_{t}$ of the boundary modes commutes with the descent at all times:
\[
\alpha_{t}\circ\mathcal{E} \;=\; \mathcal{E}\circ\alpha_{t} \quad \text{for all } t, \qquad \mathcal{E} = E_{1}\circ\cdots\circ E_{n}, \quad E_{k}(x) = N_{k}xN_{k} + (1 - N_{k})x(1 - N_{k}),
\]
$\mathcal{E}$ the descent of the window (Proposition \ref{c:singularity.descent-what-it-erases}).

\emph{The version 1 wording, equivalent.} By Theorem \emph{Boundary-mode locality is covariance under the descent}, on every finite window this holds if and only if the dynamics is generated by elements of the abelian algebra $\mathcal{D}$ that the $\sigma_{k} = 2N_{k} - 1$ generate, diagonal in the joint number basis: the internal dynamics is generated by the geometric data the framework supplies on $E$, the boost bivector field $x \mapsto \sigma_{r}(x)$ and its products. Cross-point products such as $\sigma_{j}\sigma_{k}$ correlate phases without moving occupation; hopping $v_{j}^{*}v_{k} + v_{k}^{*}v_{j}$, pairing $v_{j}^{*}v_{k}^{*} + v_{k}v_{j}$ and every other generator with a component outside $\mathcal{D}$ are excluded by the one clause. Within the class of generators the framework supplies, the exclusion of hopping is a theorem (Theorem \emph{The class-relative no-go: no supplied generator hops} [T5, Thm C.8]); what the premise adds is the exclusion of unsupplied generators, pairing terms and the possible differential dynamics, which Proposition \emph{(L$_E$) as the stabiliser and the tangent} [T5, Prop 13.4] places in the tangent of the vacuum orbit and in its stabiliser.

\emph{What it forbids.} By Theorem \emph{Number-diagonal dynamics has no exponential operator spreading}, every dynamics the premise admits carries an evolved mode operator only to the modes directly coupled to it, and every local commutator and out-of-time-order correlator is almost periodic and bounded: the premise forbids exponential real-time commutator growth. The real-time spreading law $t_{*} = \kappa^{-1}\ln N$ of Proposition \emph{The spreading law, and the chaos bound met, not saturated} [T8, Asm 6.1] is therefore not realised under it, and the scrambling time is a modular-time statement: the growth $e^{\kappa t}$ is a rate of the rotor flow, the KMS-strip continuation of the phase flow that is the windows' time evolution (Theorem \emph{The phase flow is the time evolution; the rotor flow is its KMS continuation}). The premise is kept over fast real-time scrambling; the two cannot both hold.
\end{postulate}

\begin{postulate}[(P): univalence]\label{c:singularity.p-univalence}
Physical states of the boundary modes are even: $\omega \circ \alpha = \omega$, equivalently $\omega$ annihilates the odd sector $\Aod$.
\end{postulate}

\begin{postulate}[(R): regularity at the fixed locus]\label{c:singularity.r-regularity-fixed-locus}
For every mode $k$ and every pair $a, b$ of elements of its ladder $\{v_k, v_k^{*}\}$, the correlation function $F_{a,b}(t) := \omega\bigl(a\,\alpha_t(b)\bigr)$ of the Killing flow $\alpha_t = \operatorname{Ad}(e^{it\omega_k K_k})$, an entire function of $t$, continues in imaginary Killing time to a single-valued function on the regular Euclidean section of the Schwarzschild gravitational field on flat spacetime (Proposition \emph{The Euclidean section, and its period} [T5, Prop 7.18]; Definition \emph{Gauge Theory Gravity: the gravitational field on flat spacetime}): its two orderings are the two sides of the cut, and they meet after one period,
\[
F_{a,b}\bigl(t + i\,p\bigr) \;=\; \omega\bigl(\alpha_t(b)\,a\bigr), \qquad p = \frac{2\pi}{\kappa} .
\]
In words: each single-mode marginal is KMS for the Killing flow at exactly the period the Euclidean section's smoothness at its tip dictates. The regularity is local, posed at the tip of the Euclidean section (a polar chart about the tip of the Wick-rotated flat background); for a state that is KMS for the Killing flow, this is also the period at which it is regular across the horizon locus in the Painlev\'e-Gullstrand chart.

\emph{Note (continuum restatement).} Posed on the continuum of the horizon's null generators (Definition \emph{Continuum mode space} [T7, Def 4.1]), the content of (R) for the eternal (Kruskal) hole follows from Premise \emph{Premise (N$_{+}$): the null ground state on horizon generators}: by Theorem \emph{Positive null energy gives the (R) premises}(iii) the generators' modes are then jointly KMS for the Killing flow at $\beta = 8\pi M$. The wording above, on the ladders $\{v_{k}, v_{k}^{*}\}$ with a phase-acting flow, is a discretisation of that statement, recovered only for modes of vanishing Mellin width; (N$_{+}$) itself cannot be posed on point spectrum. A hole formed by collapse is not covered by (N$_{+}$), and (R) remains a premise. (The whole generator of the eternal horizon, which that case of (N$_{+}$) uses, lies beyond one flat chart; see that premise.)
\end{postulate}

\section{The descent is the parity map}

\phantomsection\label{c:note.state-selection.setting-descent-parity-map}
The first observation is elementary and, as far as the chain is concerned, structural: the descent installed in Step 4a is not merely \emph{a} projection that happens to kill coherences. It is the canonical $\mathbb{Z}_2$ average of the algebra's own grading.

\begin{remark}[What this says about the two ledgers]\label{c:note.state-selection.what-says-about-two-ledgers}
The chain's slogan is that the metric level keeps the occupation records and loses the coherences. Theorem \ref{c:singularity.descent-parity-average} says something sharper and more structural: \emph{the metric level is the parity even level}, and the ladder sector that Rule 1 adjoined is exactly the odd sector. The one-way character of the descent \cite[Fact 5.9]{ButlerT5} is then not a separate fact to be imported but the statement that a $\mathbb{Z}_2$ average cannot be undone: the even subalgebra retains no trace of the odd part it averaged away. It also explains, at a glance, why the descent is the \emph{unique} trace-compatible conditional expectation onto the commutant of $P$: it is the unique one compatible with the grading, and the grading is not a choice but the Clifford algebra's own.
\end{remark}

\section{The form is forced; only one number is free}

\phantomsection\label{c:note.state-selection.setting-form-forced-only-one-number}
We can now dispose of the question that Step 4d's hypothesis (ii) appears to answer. It does not need answering.

\begin{remark}[The true reading, and why it sharpens the gap]\label{c:note.state-selection.true-reading-why-it-sharpens-gap}
Theorem \ref{c:singularity.descended-form-state-independent} should be read as deflationary, not triumphant. For a single fermionic mode, ``thermal'' is not a discriminating property: the state space of one mode, after the descent that the chain performs anyway, is one-dimensional, and every point of it is a Gibbs state. It follows that no part of the physical content of Step 4d lies in the \emph{shape} of the reduced density matrix. All of it lies in the number $\beta_k$, in the claim that $\beta_k$ is the same for all modes, and in the claim that its value is $2\pi/\kappa$. The gap is therefore not ``which class of states'', as the status line ``up to a state'' suggests; it is ``which point of a line, and why the same point at every mode''. Everything that follows is about that line.

This is a sharpening of the gap and not a closing of it, and it is worth being explicit that the sharpening cuts against a certain kind of rhetorical comfort: one cannot claim the Fermi-Dirac \emph{form} as a prediction of the chain, because the descent manufactures that form from any input whatsoever. What the chain can claim is the temperature, and only under the premises named below.
\end{remark}

\subsection*{Univalence, and the marginals upstairs}
\phantomsection\label{c:note.state-selection.setting-univalence-marginals-upstairs}
Theorem \ref{c:singularity.descended-form-state-independent} concerns the descended marginal. Step 4d wants a statement about the radiation itself, that is about the ascended marginal, and for that one premise is needed. It is the weakest of the standard ones.

\phantomsection\label{c:note.state-selection.setting-univalence-marginals-upstairs-2}
Premise \ref{c:singularity.p-univalence} is available to the chain for a reason the chain itself supplies: the ledger is a CAR algebra by the tower limit \cite[Thm.~5.6]{ButlerT3}, its grading is the Clifford grading of Definition \ref{c:obj.singularity.grading-parity-automorphism}, and by Theorem \ref{c:singularity.descent-parity-average} its even part is exactly the metric-level algebra. What the univalence rule asserts is that the odd sector carries no expectation values, and Theorem \ref{c:singularity.descent-parity-average} says that this is the same as asserting that the ascended and descended readings agree \emph{per mode}. That is a considerably weaker demand than boost invariance of the whole state, and it is not a choice of state: it is a statement about which elements of the algebra are observables.

\begin{remark}[What was bought]\label{c:note.state-selection.what-was-bought}
Theorem \ref{c:singularity.ascended-marginals-are-diagonal} replaces hypothesis (ii), a symmetry assumption about the state, with (P), a statement about which elements of the ledger are observable. It is strictly weaker: (ii) constrains the whole state, (P) constrains one complex number per mode. It is also more robust, and this matters for the chain: (ii) is an equilibrium condition, and the evaporating hole is not in equilibrium, while (P) holds throughout an arbitrary non-stationary history. the section ``What equilibrium cannot supply: correlations'' shows that this robustness is not a luxury but a necessity, because the exact equilibrium condition is provably too strong for what Steps 5 and 6b need.
\end{remark}

\section{One temperature, and the rigidity of the number}

\phantomsection\label{c:note.state-selection.setting-one-temperature-rigidity-number}
Theorem \ref{c:singularity.ascended-marginals-are-diagonal} leaves a function $k \mapsto \beta_k$. Two questions remain: why is it constant, and what fixes the constant.

\begin{fact}[Imported: the general statement]\label{c:state-selection.imported-general-statement}
For a $C^{*}$-dynamical system the $(\alpha,\beta)$-KMS states at fixed $\beta$ form a weak$^{*}$-compact convex simplex whose extreme points are exactly those whose GNS von Neumann algebra is a factor, distinct extremal KMS states being disjoint; in particular a factor state is extremal, and uniqueness of the $\beta$-KMS state is the statement that the simplex is a point. \emph{Provenance:} Bratteli and Robinson \cite[Thm.~5.3.30]{BR2}; the KMS condition itself is Haag, Hugenholtz and Winnink \cite{HHW1967}. The companion statement that two faithful normal states have the same modular automorphism group only if they differ by an element affiliated to the centre is Connes' Radon-Nikodym cocycle theorem \cite[Thm.~1.2.1]{Connes1973}, which for a factor forces the two states to be equal.
\emph{Source:} O. Bratteli and D. W. Robinson, Operator Algebras and Quantum Statistical Mechanics 2: Equilibrium States. Models in Quantum Statistical Mechanics, 2nd ed., Springer (1997)., Thm.~5.3.30; R. Haag, N. M. Hugenholtz and M. Winnink, On the equilibrium states in quantum statistical mechanics, Commun. Math. Phys. 5 (1967), 215--236.; A. Connes, Une classification des facteurs de type III, Ann. Sci. 'Ec. Norm. Sup'er. (4) 6 (1973), 133--252., Thm.~1.2.1.
\end{fact}

\section{Fixing the rate: one premise, and its habitat inside the framework}

\phantomsection\label{c:note.state-selection.setting-fixing-rate-one-premise-habitat}
The calibration converts rapidity to time. What remains is to fix the modular period in rapidity, and here the input is a regularity condition. This section states it, shows that the framework supplies its habitat rather than importing it wholesale, and derives the temperature.

\begin{proposition}[The Euclidean partner is curvature selected]\label{c:state-selection.euclidean-partner-curvature-selected}
Let $\sigma_r$ be the unit boost bivector of the radial plane and $I$ the pseudoscalar. Then:
\begin{enumerate}
\item $(I\sigma_r)^{2} = -1$, and $\{\sigma_r,\, I\sigma_r\}$ is precisely the doubled eigenplane pair of the shape operator, $S(B) = -2B$ on it, of \cite[Thm.~6.1]{ButlerT5};
\item the rotor $R_\theta := \exp(\theta\, I\sigma_r/2) = \cos(\theta/2) + I\sigma_r \sin(\theta/2)$ satisfies $R_{2\pi} = -1$ and $R_{4\pi} = +1$, so the rotor has period $4\pi$ while its conjugation action $x \mapsto R_\theta x R_\theta^{-1}$ has period $2\pi$;
\item $R_\theta$ is the analytic continuation of the boost rotor $\exp(\lambda\sigma_r/2)$ under $\lambda \mapsto i\theta$ in the sense that both are one-parameter rotor flows in the same eigenplane pair, the hyperbolic one generated by the member with square $+1$ and the elliptic one by the member with square $-1$.
\item On spinors, each chiral half: with the unit of the continuation $\lambda \mapsto i\theta$ read as the conformal spinors' unit $I_{6} = I_{4}E$ (with $I_{4} = I$ and $E = e_{+}e_{-}$ the bivector of the added plane), the continued boost acts on the two chiral halves of the conformal spinors as $R_\theta$ and $R_{-\theta}$, the sense fixed by the chirality (Proposition \emph{The Wick-rotated boost as the angular rotor on each chiral half}(ii)).
\end{enumerate}
In its action on spinors, on each chiral half, the Euclidean rotation whose angle carries the periodicity is therefore not an imported device: it is the second member of the pair that Step 4b's curvature equation already selects, and it is distinguished from the boost by exactly the sign of its square, the same criterion \cite[Cor.~6.4]{ButlerT5} uses to pick out $\sigma_r$. On vectors, and for every holonomy of a real connection, it is not: there $R_\theta$ rotates along the bifurcation sphere, not around it, and the Euclidean rotation around the sphere is realised by no real rotor of $\Cl(1,3)$ (Proposition \ref{c:singularity.euclidean-rotor-where-it-lies}(3)).
\end{proposition}
\begin{proof}
(1) $I$ commutes with every even element and $I^{2} = -1$ in $\Cl(1,3)$, so $(I\sigma_r)^{2} = I^{2}\sigma_r^{2} = -1$ using $\sigma_r^{2} = +1$. The eigenplane statement is \cite[Thm.~6.1]{ButlerT5} (\ref{c:singularity.curvature-gauge-field-exact-shape-decoupling}), whose spectrum is $(-2,-2,1,1,1,1)$ with the doubled eigenvalue attained exactly on $\operatorname{span}\{\sigma_r, I\sigma_r\}$. (2) With $(I\sigma_r)^{2} = -1$ the exponential series splits into the cosine and sine series, giving the displayed rotor; $R_{2\pi} = \cos\pi + I\sigma_r\sin\pi = -1$, and conjugation by $-1$ is the identity, so the conjugation action closes at $2\pi$ while the rotor itself closes at $4\pi$. (3) The two flows are the two real one-parameter subgroups of the rotor group of the plane pair, related by the substitution $\lambda = i\theta$ in the exponential; the substitution exchanges $\cosh, \sinh$ with $\cos, \sin$ exactly because the two generators differ by the factor $I$ with $I^{2} = -1$. (4) Read with the conformal spinors' unit $I_{6} = I_{4}E$ as the unit of the continuation, the continued boost is $\exp(\theta I_{6}\sigma_r/2) = \tfrac12(1+E)R_\theta + \tfrac12(1-E)R_{-\theta}$, which is Proposition \emph{The Wick-rotated boost as the angular rotor on each chiral half}(ii): on the chiral half $E = +1$ it is $R_\theta$ and on $E = -1$ it is $R_{-\theta}$, so on spinors the Euclidean rotation is carried by the second member, with the sense fixed by the chirality. The restriction to spinors is forced: by Proposition \ref{c:singularity.euclidean-rotor-where-it-lies}(3) the rotor $R_\theta$ fixes $e_0$, $e_1$ and rotates $e_2$ into $e_3$, along the bifurcation sphere, and no real rotor of $\Cl(1,3)$ realises the rotation around it, while the continued element is no rotor on vectors (Proposition \emph{The Wick-rotated boost as the angular rotor on each chiral half}(iii)).

\end{proof}

\begin{theorem}[The rate under the geometric modular condition, with hypothesis (iii) presupposed]\label{c:state-selection.rate-discharge-hypothesis-iii}
Assume (P) and (R), the latter in T05's form (\ref{c:singularity.r-regularity-fixed-locus}: on each mode the two orderings of every ladder correlator meet after one period of the regular Euclidean section), and that the state satisfies the GMC of Definition \ref{c:obj.singularity.geometric-modular-condition-g}. Then the modular period in rapidity is exactly $2\pi$, so by Corollary \ref{c:singularity.universality-temperature-under-g} and the calibration of \cite[Prop.~6.2]{ButlerT5},
\[
\beta \;=\; \frac{2\pi}{\kappa} \;=\; 8\pi M , \qquad T \;=\; \frac{\kappa}{2\pi} \;=\; \frac{1}{8\pi M} \;=\; \TH ,
\]
Hypothesis (iii) of \cite[Thm.~7.3]{ButlerT5}, that the state be faithful, that is cyclic and separating for the one-sided algebra, is presupposed here and not derived: a state has a modular automorphism group only if it is faithful, so the GMC contains it. Its discharge without the GMC is T05's Theorem \emph{The rate, and the discharge of hypothesis (iii)} [T5, Thm 7.20], from T05's (R) on each mode.
\end{theorem}
\begin{proof}
By the GMC the state is KMS for the boost's phase flow $\operatorname{Ad}(e^{itK})$ at some $\beta$, one for all modes (\ref{c:singularity.universality-temperature-under-g}), and by Theorem \ref{c:singularity.kms-uniqueness-factor} it is the corresponding Gibbs state; the KMS condition says its correlation functions are analytic in the strip of width $\beta$ and periodic, with the graded sign (the state being even by (P), \ref{c:singularity.p-univalence}), under translation by $i\beta$ in the flow parameter. Under the continuation $\lambda \mapsto i\theta$ of Proposition \ref{c:state-selection.euclidean-partner-curvature-selected}(3), read on spinors as in its clause (4), that translation acts on each chiral half as the rotor of the second member of the curvature-selected pair, by the angle $\theta = \beta$, in one sense on one half and the opposite sense on the other; on vectors it is no real rotation, so the period is read in this action on spinors. By Proposition \ref{c:state-selection.euclidean-partner-curvature-selected}(2) the conjugation action of the rotation closes at $2\pi$ and at no smaller angle, so a continuation that is single valued and regular at the fixed locus of the flow, which is premise (R) in the form T05 states it (\ref{c:singularity.r-regularity-fixed-locus}: on each mode the two orderings of every ladder correlator are glued after one period of the regular Euclidean section), forces $\beta = 2\pi$ in rapidity: a value $\beta \neq 2\pi$ makes the continued correlation function multivalued around the fixed locus, equivalently gives the Euclidean section a conical excess or deficit there. The conversion to Killing time is \cite[Prop.~6.2]{ButlerT5} (\ref{c:singularity.plane-rate}). Faithfulness: the GMC presupposes it, since a state has a modular automorphism group only if it is faithful, so hypothesis (iii) is an input of this argument and not an output. That $\beta = 2\pi$ is finite and nonzero, so that the Gibbs matrix has full rank, is consistent with it but derives nothing, the argument having started from the modular group. Hypothesis (iii) is discharged without the GMC by Theorem \emph{The rate, and the discharge of hypothesis (iii)} [T5, Thm 7.20], from T05's (R) on each mode. The note's own wording of (R), regularity of the continued two-point functions at the fixed locus, is satisfiable and contentful on the boundary-mode algebra only when the two orderings are glued, and then on the whole window it adds vanishing inter-mode two-point functions, which this argument does not use (\emph{The folded note's (R) and (L$_E$) against T05's}, (i), (ii) and (iv)); T05's (R) is the premise assumed here.

\end{proof}

\begin{corollary}[Where the $+1$ of Fermi-Dirac comes from]\label{c:state-selection.where-fermi-dirac-comes}
The rotor implementing one Euclidean period is $R_{2\pi} = -1$, which acts as $+1$ on the even sector $\Aev$ and as $-1$ on the odd sector $\Aod$. The ladder's own generators are odd, so their correlation functions are antiperiodic around the modular period while the metric-level observables are periodic. This is the structural origin of the Fermi-Dirac denominator $e^{\omega/T} + 1$ rather than the Bose-Einstein $e^{\omega/T} - 1$, and hence of the first wrinkle recorded in \cite[Rem.~7.5]{ButlerT5}: the statistics of the directly produced spectrum is fixed by the double cover, and the bosonic sector has to be reached by bosonisation, as Step 4e does.
\end{corollary}
\begin{proof}
The value $R_{2\pi} = -1$ is Proposition \ref{c:state-selection.euclidean-partner-curvature-selected}(2). Conjugation by $-1$ is trivial, so the even sector, which is where the metric-level observables live by Theorem \ref{c:singularity.descent-parity-average}, is periodic. The odd sector is acted on by the rotor in the spinor module, where $-1$ acts as $-1$; the antiperiodicity of odd correlation functions around the modular period is the standard consequence, and the Matsubara reading of that antiperiodicity is the Fermi-Dirac factor. The last step is the standard reading and is flagged as such rather than proved here.

\end{proof}

\begin{remark}[The true price of (R)]\label{c:note.state-selection.true-price-r}
Premise (R) is not derived from the transport principle, and this note does not claim otherwise. It is the same premise that every derivation of the Hawking temperature uses: regularity of the Euclidean section at the bifurcation point in \cite{GibbonsHawking1977, HartleHawking1976}, regularity on the horizon in \cite{Unruh1976}, nonsingularity of the quasi-free state in \cite{KayWald1991}, the short-distance condition in \cite{FredenhagenHaag1990}. Dropping it gives the Boulware alternative \cite{Boulware1975}, which in this framework is the $\beta = \infty$ endpoint of the line: an unoccupied ledger, a descended state that is pure, a deficit of zero by \cite[Thm.~9.3]{ButlerT5}, and no flux. What the framework contributes is not the premise but its habitat. In a metric formulation the Euclidean section is a separate construction; here, by Proposition \ref{c:state-selection.euclidean-partner-curvature-selected}, the object whose periodicity is being demanded is the second member of the very eigenplane pair that Step 4b's curvature equation selects, and its period is fixed by the double cover of the rotor group, not by a coordinate choice. The premise is thereby stated in the same algebra as the rest of the chain, and its content is exactly one number.
\end{remark}

\begin{remark}[The dictionary]\label{c:note.state-selection.dictionary}
The standard alternatives sit on the line as follows. $\beta_k = \infty$ for all $k$: the Boulware analogue, ladder unoccupied, no radiation, zero deficit, excluded by (R). $\beta_k = 2\pi/\kappa$ for all $k$: the Hartle-Hawking and Unruh analogue, the GMC state, selected by (R). $\beta_k$ finite but not constant in $k$: no geometric modular flow by Corollary \ref{c:singularity.universality-temperature-under-g}, so a state for which the modular clock is not the Killing boost at all. $\beta_k = 0$: the tracial state, infinite temperature, excluded with the other boundary case. The chain's contribution to the dictionary is that the entries are points of an explicitly parametrised line rather than incommensurable constructions.
\end{remark}

\section{Purity is an output, not an input}

\phantomsection\label{c:note.state-selection.setting-purity-output-not-input}
Hypothesis (i) of \cite[Thm.~7.3]{ButlerT5} asks for a pure state on the doubled system. It need not be asked for.

\begin{theorem}[The partner and the purification are canonical]\label{c:state-selection.partner-purification-are-canonical}
Let $\omega$ be a faithful state on the mode factor $\Aa \cong \Mat_2(\C)$, with density matrix $\rho = \operatorname{diag}(\lambda_0^{2}, \lambda_1^{2})$, $\lambda_i > 0$. Then in the GNS representation of $\omega$:
\begin{enumerate}
\item the representation space is $\C^{2}\otimes\C^{2}$, the cyclic vector is $\Omega = \sum_i \lambda_i |ii\rangle$, and $\Omega$ is cyclic and separating for $\pi(\Aa) = \Aa \otimes 1$;
\item the commutant $\pi(\Aa)' = 1 \otimes \Aa$ is the partner copy, and is generated from $\pi(\Aa)$ by the modular conjugation, $\pi(\Aa)' = J\pi(\Aa)J$;
\item any two purifications of $\omega$ are related by a unitary acting on the purifying factor alone, so $\Omega$ is unique up to a unitary in the commutant;
\item when $\omega$ is the Gibbs state at $\beta$, $\Omega$ is the thermofield double, $\lambda_0 = \cos\theta$, $\lambda_1 = \sin\theta$ with $\tan\theta = e^{-\omega_k/2T}$, which is exactly the state produced by \cite[Lem.~7.2]{ButlerT5}.
\end{enumerate}
Consequently the global purity of hypothesis (i), the existence of the partner, and the thermofield form of the global state are consequences of the one-sided state, not independent choices; and the ``two ends of the pair'' of Step 4d are the algebra and its commutant in the state's own GNS representation.
\end{theorem}
\begin{proof}
(1) and (2) are \cite[Lem.~7.1]{ButlerT5} (\ref{c:singularity.modular-theory-matrix-factor-proved-outright}) read in the direction of construction rather than of analysis: that lemma computes $\Delta = \rho\otimes\bar\rho^{-1}$ and $J|ij\rangle = |ji\rangle$ for exactly this $\Omega$, and establishes that $\Omega$ is cyclic and separating; the identification $\pi(\Aa)' = J\pi(\Aa)J$ is Tomita's theorem, which in this finite-dimensional case is the content of that lemma's polar decomposition. (3) If $|\Psi\rangle \in \C^2\otimes\mathcal{K}$ purifies $\rho$ then its Schmidt coefficients are the $\lambda_i$, so $|\Psi\rangle = (1\otimes W)\Omega$ for an isometry $W$ carrying one Schmidt basis to the other; two purifications on the same purifying space differ by a unitary there. (4) Substituting $\rho = \operatorname{diag}(1, e^{-\omega_k/T})/Z$ gives $\lambda_1/\lambda_0 = e^{-\omega_k/2T} = \tan\theta$, which is the state of \cite[Lem.~7.2]{ButlerT5}.

\end{proof}

\begin{remark}[Israel's reading, upgraded]\label{c:note.state-selection.israels-reading-upgraded}
\cite[Rem.~7.4]{ButlerT5} records that the natural state satisfying (i) to (iii) is the thermofield double, Israel's reading of the eternal black hole \cite{Israel1976}. Theorem \ref{c:state-selection.partner-purification-are-canonical} makes the word ``natural'' precise: the thermofield double is not one state among those satisfying the hypotheses, it is the canonical purification of the one-sided state, and the partner is not an extra system but the commutant. In the chain's own terms, the adjoined ladder's ``other end'' is produced by the modular conjugation of the state, so the ascent supplies the algebra and the state supplies the partner.
\end{remark}

\section{What equilibrium cannot supply: correlations}

\phantomsection\label{c:note.state-selection.setting-what-equilibrium-cannot-supply-correlations}
There is a result in the other direction, and it is the most consequential of the note for the chain as a whole. The temptation, once the GMC is available, is to make it the whole state-selection principle. That would be a mistake, and provably so.

\begin{remark}[Why this rescues rather than damages the chain]\label{c:note.state-selection.why-rescues-rather-than-damages-chain}
Read carelessly, Theorem \ref{c:singularity.exact-equilibrium-over-determines-boundary-modes} looks like an obstruction: the state that thermality wants is the state that has no information in it. Read carefully, it is the reason the architecture of this note matters. Had per-mode thermality been derived from the exact equilibrium condition, as hypothesis (ii) suggests, then thermality and correlation would have been in direct competition and the chain could not have had both. Theorem \ref{c:singularity.ascended-marginals-are-diagonal} shows that per-mode thermality does not rest on equilibrium at all: it rests on parity, which survives an arbitrary non-stationary history. The correct state class for the chain is therefore
\[
\bigl\{\ \text{even states, per-mode populations at } \beta = 2\pi/\kappa, \ \text{correlations free}\ \bigr\},
\]
a class in which every single quantum is exactly thermal and the joint state is free to carry everything the evaporation writes. That is exactly the class the finite-ledger model of \cite[\S 12]{ButlerT5} inhabits, and \cite[Lem.~12.1]{ButlerT5} is the model's proof that it does.
\end{remark}

\subsection*{Why the populations do not drift}
\phantomsection\label{c:note.state-selection.setting-why-populations-do-not-drift}
One loose end remains. Premise (P) fixes the \emph{form} of the marginals for all time; the GMC fixes the \emph{populations} but only in equilibrium. In a non-stationary history the populations are set when the mode is created, by the local geometry, which is the injection input (I) of \cite[\S 11.3]{ButlerT5}. What keeps them from drifting afterwards?

\begin{remark}[The status of (L$_E$)]\label{c:note.state-selection.status-l-e}
(L$_E$) is a premise, not a theorem, and it is the true location of a genuine assumption: it says that the framework supplies no hopping generator such as $v_j^{*}v_k$ between distinct points of $E$. That is a statement about what the localisation theorems provide, and they do provide the ladder pair pointwise; it is not a proof that no such coupling can arise from a dynamics not yet written down. Recorded as a programme item in the section ``What remains open''. What Proposition \ref{c:singularity.number-diagonal-dynamics-models-scrambling-derived} buys, if (L$_E$) holds, is worth the exposure: it converts the model's most questionable-looking choice, that the scrambling is phase-type, into the general case, and with it upgrades \cite[Lem.~12.1]{ButlerT5} from a statement about a particular circuit to a statement about the ledger.
\end{remark}

\section{The gap, relocated}

\phantomsection\label{c:note.state-selection.setting-gap-relocated}
Collecting: Step 4d's three hypotheses and the temperature now stand as in the table accompanying \ref{c:singularity.number-diagonal-dynamics-models-scrambling-derived}.

 \begin{tabular}{l>{ \arraybackslash}p{4.9cm}>{ \arraybackslash}p{6.3cm}}
\toprule
Item in \cite[Thm.~7.3]{ButlerT5} & Was & Now \\
\midrule
(i) purity & assumed & the canonical GNS purification of the one-sided state; the partner is the commutant (Thm \ref{c:state-selection.partner-purification-are-canonical}) \\
(ii) boost invariance & assumed & not needed for the form at all (Thm \ref{c:singularity.descended-form-state-independent}); the per-mode diagonality it supplied follows from univalence (Thm \ref{c:singularity.ascended-marginals-are-diagonal}); and in exact form it is too strong (Thm \ref{c:singularity.exact-equilibrium-over-determines-boundary-modes}) \\
(iii) cyclic and separating & assumed & equivalent to $0 < \beta < \infty$, hence a consequence of (R) (Thm \ref{c:state-selection.rate-discharge-hypothesis-iii}) \\
the value $\TH$ & fixed by the calibration & cannot be fixed algebraically (Thm \ref{c:singularity.rate-rigidity-temperature-not-algebra-intrinsic}); fixed by (R) plus the calibration, uniquely (Thm \ref{c:state-selection.rate-discharge-hypothesis-iii}) \\
correlations & implicit, supplied by Step 6b's axioms & provably not supplied by any equilibrium condition (Thm \ref{c:singularity.exact-equilibrium-over-determines-boundary-modes}); dynamical input, priced where it is used \\
\bottomrule
\end{tabular}

\emph{Step 4d after this note. The status line ``up to a state'' becomes ``proved from (P) and (R)''.}

The proposed status line for the chain table of \cite[Table 1]{ButlerT5} is therefore \emph{proved from (P), (R)} rather than \emph{up to a state}, with the two premises named in the row's own text, in exactly the manner Steps 6a and 6b already name (C1) to (C3) and (U), (C), (F). That is not a cosmetic change: it moves Step 4d from the class of results that await a state to the class of results that hold given premises the reader can inspect, which is the standard the rest of the chain is held to.

What has been relocated rather than removed is this. The chain's residual state dependence is now entirely in the correlations, and the correlations are the business of Step 6b, where the rate model already prices them by (U), (C) and (I). Theorem \ref{c:singularity.exact-equilibrium-over-determines-boundary-modes} adds one thing to that pricing which was not previously visible: the correlations cannot be obtained from any equilibrium principle whatsoever, so no sharpening of Step 4d will ever supply them, and the search for a state-selection principle that delivers both thermality and the Page curve is a search for something that provably does not exist in the equilibrium category. Steps 4d and 6b are not two instalments of one problem; they are two different problems, and this note closes the first at the price of showing that the second cannot be closed the same way.

\section{What remains open}

\phantomsection\label{c:note.state-selection.setting-what-remains-open}
\begin{enumerate}
\item \emph{Premise (R) is not derived, and cannot be, from the singular side.} It is the standard premise, stated in the framework's own algebra (Proposition \ref{c:state-selection.euclidean-partner-curvature-selected}), but the transport principle does not entail it, and the reason is structural rather than a shortcoming of the present proof. The only finite rate the geometry supplies to the boost flow is the surface gravity, a horizon quantity: on the singular side the connection rate diverges as $\sqrt{M/2r^{3}}$ and offers no finite period to identify with $2\pi$. The locus at which (R) has purchase is the fixed locus of the boost flow, the bifurcation surface at $r = 2M$, and the Euclidean section on which the periodicity argument runs is the region $r \geq 2M$, which does not contain the singular boundary at all. There, by the chain's own Step 1, the gauge field is invertible and unimodular and division does not fail. A theory of the loci where division fails is therefore silent, by its own diagnosis, at the locus where the temperature is calibrated. In particular, strengthening the rain-finiteness of the chain's null dyad \cite[Prop.~8.23, Rem.~8.24]{ButlerT4} from a statement about the dyad to a statement about admissible states at $E$, which is the question \cite{ButlerT4} itself poses and a reasonable item for the geometry paper (T6), would discharge faithfulness and exclude the endpoints $\beta \to 0$ and $\beta \to \infty$, but would not fix the value: it is a bound, not a calibration.

The live route to a better-placed premise runs through Step 3's landing zone instead. Bisognano-Wichmann does not assume the factor $2\pi$: it derives it, from the spectrum condition together with covariance and locality \cite{BW1976, Sewell1982}. The ascent lands in $\Cl(2,4)$, whose parabolic stabiliser of the adjoined null direction contains exactly the translations such a spectrum condition needs, so the premise could be restated as positivity of the ledger's energy with respect to a generator the ascent itself installed. That would still be an import, positivity of energy being no Clifford fact, but it would be stated in the algebra the chain produces rather than in the metric geometry. Transcribing the theorem needs an action of the relevant subgroup on the ledger, a wedge and locality structure, and analyticity of the state, each of which is real work and any of which could fail; the natural home is the twistor paper, not T6.
\item \emph{Premise (L$_E$) is a statement about available generators.} See Remark \ref{c:note.state-selection.status-l-e}. A no-go excluding inter-point hopping generators on $E$, in the manner of the series' other no-go results, would settle it.
\item \emph{The infinite ledger.} Theorems \ref{c:singularity.descent-parity-average} to \ref{c:singularity.ascended-marginals-are-diagonal} and \ref{c:singularity.exact-equilibrium-over-determines-boundary-modes} are stated for the CAR algebra and their proofs do not use finiteness. Theorems \ref{c:singularity.kms-uniqueness-factor}, \ref{c:singularity.rate-rigidity-temperature-not-algebra-intrinsic} and \ref{c:state-selection.partner-purification-are-canonical} are proved here in the matrix factor, which is the setting Step 4d works in; the general versions are the imported Facts \ref{c:state-selection.imported-general-statement}, and in the genuine field-theoretic setting the horizon algebra is a type III$_1$ factor, where the Gibbs form is unavailable and the KMS condition must be used directly. Transcribing the section ``One temperature, and the rigidity of the number'' and the section ``Fixing the rate: one premise, and its habitat inside the framework'' into that setting is the natural next technical step, and \cite{KayWald1991} is the template.
\item \emph{Universality across the sphere.} Corollary \ref{c:singularity.universality-temperature-under-g} gives one $\beta$ for the whole ledger given the GMC, but the ledger is indexed by $E \cong \R \times S^{2}$ and the argument treats the sphere directions as spectators. That is correct for Schwarzschild by symmetry; for Kerr the boost weights vary over the horizon and the corresponding statement is the rigidity of the temperature over a rotating horizon, which the chain has not touched.
\item \emph{The greybody interface.} Nothing here bears on the partial-wave reconstruction flagged in \cite[Rem.~8.3]{ButlerT5}; it remains assembly level.
\end{enumerate}

\section{Machine verification}

\phantomsection\label{c:note.state-selection.setting-machine-verification}
Every computational identity in this note has been checked in exact rational arithmetic, in the manner of the scripts deposited with the series' papers, by \texttt{verify\_state\_selection.py} \cite{ButlerVerifySS}, deposited alongside and using the series' shared library \texttt{GA.py}. Forty-one checks, no floating point, nonzero exit on any failure. Where an exponential occurs it is carried as an exact rational variable, $x = e^{-\beta\omega} \in \mathbb{Q}$, and never evaluated numerically.

Coverage. Theorem \ref{c:singularity.descent-parity-average}: the identity $E(x) = \tfrac12(x + \sigma x \sigma)$ on a generic rational matrix; idempotence; the conditional-expectation property $E(axb) = aE(x)b$ for even $a, b$; the range computed as the even span through $E(1) = 1$ and $E(\sigma) = \sigma$; the kernel computed as the odd span through $E(v) = E(v^{*}) = 0$ together with the exhaustive statement that $E$ kills exactly the off-diagonal part; and the hyperbolic atom's CAR with $\sigma = e_0e_1$ blade-exactly in the shared Clifford engine. Theorem \ref{c:singularity.descended-form-state-independent}: for a density matrix with nonzero coherences, $E(\rho)$ is diagonal, is a trace-one state, and equals $\operatorname{diag}(1, x)/(1+x)$ at the $x$ the theorem names, with the injectivity of $x \mapsto p$ checked on a rational witness pair. Theorem \ref{c:singularity.ascended-marginals-are-diagonal}: an even three-mode state, built as a rational mixture of two even pure states, is verified even by $\Gamma\rho\Gamma = \rho$, has all three single-mode marginals diagonal, and satisfies $\omega(v_k) = \omega(v_k^{*}) = 0$ at every mode; a parity-mixed control has marginal coherence $\tfrac{12}{25}$, so the premise is demonstrably doing work; and the model's pair state of \cite[\S 12]{ButlerT5}, controlled-$Z$ sign included, is verified even with all marginals diagonal. Theorem \ref{c:singularity.kms-uniqueness-factor}: the KMS identity checked on all sixteen matrix-unit pairs for $\rho \propto \operatorname{diag}(1, \tfrac{16}{9})$; the converse checked by exhibiting a matrix-unit pair at which the tracial state fails it; and the commutation $[\rho, e^{\beta H}] = 0$ that the identity forces. Corollary \ref{c:singularity.universality-temperature-under-g}: with weights $(1,2)$, equal inverse temperatures are the exact rational condition $x_2 = x_1^{2}$, the joint populations of the two-mode product are those of a single Gibbs state of $K_{\mathrm{geo}}$ exactly then, a witness pair is exhibited for which the condition fails, and the GMC state is verified to factorise. Theorem \ref{c:singularity.rate-rigidity-temperature-not-algebra-intrinsic}: the diagonal unitary commutes with $K$ and fixes the Gibbs state exactly; the exchange unitary satisfies $XKX = -K$ and carries the $\beta$-Gibbs state to the $(-\beta)$ one; and the exchange is verified not to commute with $K$, so it is outside the flow-preserving group. Proposition \ref{c:state-selection.euclidean-partner-curvature-selected}: $\sigma_r^{2} = +1$ and $(I\sigma_r)^{2} = -1$ blade-exactly in $\Cl(1,3)$; both members of the pair are shape-operator eigenvectors of eigenvalue $-2$; the rotor composition law $R(\tfrac35,\tfrac45)^{2} = R(-\tfrac{7}{25},\tfrac{24}{25})$; $R_{2\pi} = -1$ and $R_{4\pi} = +1$ obtained by squaring the quarter rotor; and the triviality of conjugation by $-1$ on a generic multivector. Theorem \ref{c:singularity.exact-equilibrium-over-determines-boundary-modes}: the two-mode Gibbs product is a fixed point of $\EN$ and equals the tensor product of its own marginals, against an entangled even control that fails both and whose purity strictly drops under descent. Theorem \ref{c:state-selection.partner-purification-are-canonical}: the thermofield vector purifies $\operatorname{diag}(\tfrac{9}{25}, \tfrac{16}{25})$; the angle satisfies $\tan^{2}\theta = \tfrac{16}{9} = e^{-\omega/T}$; the modular operator of \cite[Lem.~7.1]{ButlerT5} is recomputed as $\rho \otimes \bar\rho^{-1}$; and a unitary on the purifying factor alone leaves the marginal unchanged while changing the global vector. Proposition \ref{c:singularity.number-diagonal-dynamics-models-scrambling-derived}: a controlled-$Z$ circuit on the model's four-mode ledger preserves every joint-number probability and every single-mode marginal, commutes with the total parity, and does change the joint coherence, so it writes exactly where the descent cannot read; while an orthogonal, parity-preserving hopping rotation between two modes changes the marginal populations, exhibiting exactly what (L$_E$) excludes.

\section{Descent, parity and states}

\begin{definition}[The grading, and the parity automorphism]\label{c:obj.singularity.grading-parity-automorphism}
Write $\Aa$ for the algebra generated by the completed pair, $\Aa \cong \Mat_2(\C)$ as a factor. It carries the $\mathbb{Z}_2$ grading it inherits from the Clifford algebra: $\Aev := \operatorname{span}\{1, \sigma\}$ is the even part and $\Aod := \operatorname{span}\{v, v^{*}\}$ the odd part, where $\sigma := v^{*}v - vv^{*}$ is the bivector of Theorem \emph{Modular structure of the completed pair} [T5, Thm 6.3]. The \emph{parity automorphism} is $\alpha(x) := \sigma x \sigma$, well defined because $\sigma^{2} = 1$; it fixes $\Aev$ pointwise and acts by $-1$ on $\Aod$. On the full boundary mode space the parity automorphism is $\operatorname{Ad}(\Gamma)$ with $\Gamma := \prod_k \sigma_k$, the total fermion parity; the mode-wise refinement, one parity automorphism $\operatorname{Ad}(\sigma_k)$ per point of $E$, is Proposition \emph{(P$_{\mathrm{loc}}$), reduced: the metric-level reading on many modes} [T5, Prop 9.6]'s task, and on two or more modes the two generate different averages.
\end{definition}

\begin{theorem}[The descent is the parity average]\label{c:singularity.descent-parity-average}
For every $x \in \Aa$,
\[
E(x) \;=\; \tfrac12\bigl(x + \alpha(x)\bigr) \;=\; \tfrac12\bigl(x + \sigma x \sigma\bigr) .
\]
Consequently $E$ is the group average over $\{\id, \alpha\}$, hence the unique conditional expectation of $\Aa$ onto its fixed-point algebra $\Aev$; its range is $\Aev$ and its kernel is exactly $\Aod$, the adjoined ladder sector; and $E$ preserves every state $\omega$ with $\omega \circ \alpha = \omega$.
\end{theorem}
\begin{proof}
In the defining representation $v = E_{12}$, $v^{*} = E_{21}$, so $P = E_{22}$ and $\sigma = 2P - 1 = \operatorname{diag}(-1, +1)$; then
\[
\alpha\begin{pmatrix} a & b \\ c & d \end{pmatrix}
= \begin{pmatrix} -1 & 0 \\ 0 & 1\end{pmatrix}\begin{pmatrix} a & b \\ c & d \end{pmatrix}\begin{pmatrix} -1 & 0 \\ 0 & 1\end{pmatrix}
= \begin{pmatrix} a & -b \\ -c & d \end{pmatrix},
\]
so $\tfrac12(x + \alpha(x)) = \operatorname{diag}(a,d)$, which is the value of $E$ computed in Proposition \ref{c:singularity.descent-what-it-erases}. The average of a finite group acting by automorphisms is a conditional expectation onto the fixed-point algebra. It is moreover the unique conditional expectation of $\Aa$ onto $\Aev$, by bimodularity: if $F$ is any conditional expectation onto $\Aev$, then for odd $x$, since $\sigma \in \Aev$ and $\sigma x \sigma = -x$, bimodularity gives $\sigma F(x) \sigma = F(\sigma x \sigma) = -F(x)$, while $F(x) \in \Aev$ is fixed by conjugation with $\sigma$, so $F(x) = -F(x) = 0$: every conditional expectation onto $\Aev$ annihilates the $(-1)$-eigenspace of $\alpha$, which is $\Aod = \operatorname{span}\{v, v^{*}\}$, and agrees with the identity on $\Aev$, hence equals $E$ on $\Aa = \Aev \oplus \Aod$. State preservation is $\omega(E(x)) = \tfrac12(\omega(x) + \omega(\alpha(x))) = \omega(x)$. The script verifies the identity, idempotence, the conditional-expectation property $E(axb) = aE(x)b$ for even $a, b$, and the range and kernel, on generic rational data.

\end{proof}

\begin{theorem}[The descended form is state independent]\label{c:singularity.descended-form-state-independent}
Let $\varrho$ be the density matrix of \emph{any} state on the boundary mode space, and $\varrho_k$ its marginal at mode $k$. Then $E_k(\varrho_k)$ is diagonal in that mode's number basis, and therefore
\[
E_k(\varrho_k) \;=\; \frac{e^{-\beta_k \omega_k N_k}}{Z_k}
\]
for a unique $\beta_k \in [-\infty, +\infty]$, determined by $\beta_k\omega_k = \log\bigl(p_k/(1-p_k)\bigr)$ with $p_k$ the descended vacancy probability. (Stated for $\omega_k \neq 0$: at a zero-weight mode the right side is $\tfrac12\,\mathrm{id}$ for every $\beta_k$, so uniqueness always fails there, and existence fails unless $p_k = \tfrac12$.) The descended single-mode state is a Gibbs state of the boost generator $K_k$ whatever the ascended state was: no hypothesis on the state is used.
\end{theorem}
\begin{proof}
By Theorem \ref{c:singularity.descent-parity-average} the range of $E_k$ is the even subalgebra, which in the number basis is the diagonal. A diagonal density matrix on a two-dimensional space is $\operatorname{diag}(p_k, 1 - p_k)$, and $\beta_k \mapsto (1 + e^{\beta_k\omega_k})^{-1}$ is a bijection from $\R$ onto $(0,1)$, so every such matrix with $0 < p_k < 1$ is the stated Gibbs matrix for exactly one $\beta_k$; the two boundary values are the limits $\beta_k = \pm\infty$.

\end{proof}

\begin{fact}[Imported: the univalence superselection rule]\label{c:singularity.imported-univalence-superselection-rule}
In a theory whose field algebra is $\mathbb{Z}_2$ graded by fermion number, coherent superpositions of states of opposite fermion parity are not observable: the observable algebra is the even subalgebra and physical states are even. \emph{Provenance:} Wick, Wightman and Wigner \cite{WWW1952}; the graded structure of CAR algebras and their even subalgebras is \cite[\S 5.2]{BR2}. What the rule superselects is the total parity, a single $\mathbb{Z}_2$; the mode-wise strengthening that the multi-mode descent of Step 5 needs is Proposition \emph{(P$_{\mathrm{loc}}$), reduced: the metric-level reading on many modes} [T5, Prop 9.6], stated there, and is not part of this import.
\emph{Source:} G. C. Wick, A. S. Wightman and E. P. Wigner, The intrinsic parity of elementary particles, Phys. Rev. 88 (1952), 101--105.; O. Bratteli and D. W. Robinson, Operator Algebras and Quantum Statistical Mechanics 2: Equilibrium States. Models in Quantum Statistical Mechanics, 2nd ed., Springer (1997)., \S 5.2.
\end{fact}

\begin{theorem}[Ascended marginals are diagonal]\label{c:singularity.ascended-marginals-are-diagonal}
Let $\omega$ satisfy (P). Then $\omega(v_k) = \omega(v_k^{*}) = 0$ at every mode, so the ascended marginal $\varrho_k$ is diagonal in the number basis, coincides with $E_k(\varrho_k)$, and is by Theorem \ref{c:singularity.descended-form-state-independent} the Gibbs state of $K_k$ at a unique $\beta_k$ for $\omega_k \neq 0$ (at a zero-weight mode diagonality still holds, and the Gibbs form requires $p_k = \tfrac12$, Theorem \ref{c:singularity.descended-form-state-independent}'s parenthesis). Consequently partial trace and pinching coincide on every single-mode marginal, which is Proposition \emph{Pinching against partial trace: the scope of a coincidence} [T5, Prop 9.2](i) with hypothesis (ii) of Theorem \emph{Modular descent} [T5, Thm 7.3] removed; and the parity protection lemma of the finite-mode model, Lemma \emph{Parity protection} [T5, Lem 12.1], is the special case of this theorem for the model's states, which are even because every configuration in their support has even total occupation.
\end{theorem}
\begin{proof}
$v_k$ and $v_k^{*}$ are odd, so $\alpha(v_k) = -v_k$ and (P, \ref{c:singularity.p-univalence}) gives $\omega(v_k) = \omega(\alpha(v_k)) = -\omega(v_k)$, hence $\omega(v_k) = 0$, and likewise for $v_k^{*}$. The off-diagonal entries of $\varrho_k$ in the number basis are exactly $\omega(v_k)$ and $\omega(v_k^{*})$, those being the mode's two off-diagonal matrix units. Diagonality is the coincidence with $E_k(\varrho_k)$, and Theorem \ref{c:singularity.descended-form-state-independent} supplies the Gibbs form. The argument uses no finiteness and holds verbatim on the CAR algebra over $E$. For the model: its states are supported on configurations with $R_k = L_k$ at every $k$, whose total occupation $2\sum_k b_k$ is even, and a state supported in a single parity sector is even. The script verifies both clauses, with a parity-mixed control whose marginal coherence is nonzero.

\end{proof}

\begin{theorem}[KMS uniqueness at a factor]\label{c:singularity.kms-uniqueness-factor}
Let $\Aa$ be a matrix factor, $\alpha_t = \operatorname{Ad}(e^{itK})$ a flow and $\beta \neq 0$. A state $\omega = \tr(\varrho\,\cdot\,)$ is $(\alpha,\beta)$-KMS if and only if $\varrho = e^{-\beta K}/\tr e^{-\beta K}$. In particular there is exactly one KMS state for each $\beta$, and, for $K$ not a multiple of $1$, the KMS states of the flow form a line parametrised by $\beta$.
\end{theorem}
\begin{proof}
In finite dimensions the KMS condition reads $\omega(ab) = \omega\bigl(b\,\alpha_{i\beta}(a)\bigr) = \omega\bigl(b\,e^{-\beta K}a\,e^{\beta K}\bigr)$ for all $a, b$. Writing both sides as traces and using cyclicity, the left side is $\tr(b\varrho a)$ and the right side is $\tr\bigl(e^{\beta K}\varrho\, b\, e^{-\beta K} a\bigr)$, so for all $b$,
\[
b\,\bigl(\varrho\, e^{\beta K}\bigr) \;=\; \bigl(e^{\beta K}\varrho\bigr)\, b .
\]
Taking $b = 1$ gives $\varrho e^{\beta K} = e^{\beta K}\varrho$, so $\varrho$ commutes with $K$; the displayed identity then says that $\varrho\,e^{\beta K}$ commutes with everything, hence is central, hence scalar because $\Aa$ is a factor. So $\varrho \propto e^{-\beta K}$, and normalisation fixes the constant. The converse is the same computation read backwards. The script checks the identity on all sixteen matrix-unit pairs for $\varrho \propto \operatorname{diag}(1, \tfrac{16}{9})$ and exhibits a pair at which the tracial state fails it.

\end{proof}

\begin{remark}[The metric level is the parity even level]\label{c:note.singularity.metric-level-parity-even-level}
Theorem \ref{c:singularity.descent-parity-average} sharpens the chain's slogan that descent keeps the records and loses the coherences. What descent keeps is the even subalgebra, and the sector the completion installs is exactly the odd one (its ladder pair mixes the physical and adjoined planes: $k^{*} = \nO/2\nu$ against the dying direction $k$ itself, Fact \emph{Imported: one-sided completion at the singularity} [T5, Fact B.2]; this $\nu$ is T04's coupling constant, not Step 6a's cell density): the two-level split is the Clifford algebra's own $\mathbb{Z}_2$ grading, not an extra structure imposed on it. This is a single-mode statement: on the full boundary mode space the metric-level algebra the chain uses is the joint-diagonal one, the fixed points of one parity per point of $E$, strictly smaller than the total-parity-even algebra, and its licence is Proposition \emph{(P$_{\mathrm{loc}}$), reduced: the metric-level reading on many modes} [T5, Prop 9.6], not the grading alone. The one-way character recorded next as Fact \emph{Imported: the descent is a one-way conditional expectation} [T5, Fact 5.9] is then the statement that a $\mathbb{Z}_2$ average cannot be undone, and the uniqueness clause there is strengthened rather than merely imported: Fact \emph{Imported: the descent is a one-way conditional expectation} [T5, Fact 5.9]'s Takesaki clause is uniqueness among trace-preserving conditional expectations, while Theorem \ref{c:singularity.descent-parity-average} makes $E$ the unique conditional expectation onto $\Aev$ outright, trace-preserving or not; and the grading is not a choice. Step 4d turns this identification into the state-selection theorem of the section ``What selects the state''.
\end{remark}

\begin{remark}[The unconditional reading, and how it sharpens the gap]\label{c:note.singularity.unconditional-reading-how-it-sharpens-gap}
Theorem \ref{c:singularity.descended-form-state-independent} says less than it appears to, and that matters. For a single fermionic mode ``thermal'' is not a discriminating property: after the descent the mode's state space is one-dimensional and every point of it is a Gibbs state. So no part of the physical content of this step lies in the \emph{shape} of the reduced density matrix, and the chain must not claim the Fermi-Dirac form as a prediction: the descent produces that form from any input whatever. All of the content lies in the number $\beta_k$, in the claim that it is the same at every mode, and in the claim that its value is $2\pi/\kappa$. The gap is therefore not ``which class of states''; it is ``which point of a line, and why the same point at every mode''. The rest of this subsection is about that line.
\end{remark}

\begin{remark}[What hypothesis (ii) cost, and what (P) costs]\label{c:note.singularity.what-hypothesis-ii-cost-what-p}
Theorem \ref{c:singularity.ascended-marginals-are-diagonal} replaces a symmetry assumption about the state by a statement about which elements are observable. The replacement is strictly weaker: hypothesis (ii) constrains the whole state, (P) constrains one complex number per mode. It is also more robust, and Theorem \ref{c:singularity.exact-equilibrium-over-determines-boundary-modes} below shows that the robustness is necessary, not merely convenient: hypothesis (ii) is an equilibrium condition, the evaporating hole is not in equilibrium, and the exact equilibrium condition turns out to be too strong for what Steps 5 and 6b need.
\end{remark}

\section{Temperature and Hawking radiation}

\begin{definition}[The geometric modular condition (G)]\label{c:obj.singularity.geometric-modular-condition-g}
Let $K_{\mathrm{geo}} := \sum_k \omega_k K_k$ generate the Killing boost flow of the boundary modes, the flow whose bivector at each point of $E$ is the curvature-selected $\sigma_r$ of Theorem \ref{c:singularity.curvature-gauge-field-exact-shape-decoupling} and Corollary \emph{The matching: Steps 4b and 4c join by equation} [T5, Cor 6.4], with $\omega_k$ the mode's weight. A \emph{faithful} state satisfies the \emph{geometric modular condition} (G) if its modular automorphism group is the phase flow $\operatorname{Ad}(e^{itK_{\mathrm{geo}}})$ generated by $K_{\mathrm{geo}}$, not its real rotor flow, up to a single constant reparametrisation. Faithfulness is part of the condition, not a consequence of it: a state has a modular automorphism group only if it is faithful, so (G) presupposes hypothesis (iii) of Theorem \emph{Modular descent} [T5, Thm 7.3] and cannot be used to derive it.
\end{definition}

\begin{corollary}[Universality of the temperature under (G)]\label{c:singularity.universality-temperature-under-g}
A state satisfying (G) has $\beta_k = \beta$ independent of $k$ over the modes of nonzero weight: one temperature for all the boundary modes. Conversely a state with per-mode marginals $\mathrm{Gibbs}(\beta_k\omega_k)$ and no inter-mode correlation satisfies (G) only if the $\beta_k$ of the modes of nonzero weight coincide.
\end{corollary}
\begin{proof}
By Theorem \ref{c:singularity.kms-uniqueness-factor} applied to $K_{\mathrm{geo}}$, a state satisfying (G) is $\varrho \propto e^{-\beta K_{\mathrm{geo}}} = \bigotimes_k e^{-\beta\omega_k K_k}$, whose $k$-th marginal is $\mathrm{Gibbs}(\beta\omega_k)$: the same $\beta$ at every mode. Conversely the modular generator of a product state is $\log\varrho = -\sum_k \beta_k\omega_k N_k + \mathrm{const}$, proportional to $K_{\mathrm{geo}} = \sum_k \omega_k N_k + \mathrm{const}$ if and only if $\beta_k\omega_k = \lambda\omega_k$ for a single $\lambda$ and every $k$ of nonzero weight (under Premise \emph{(Q): joint quasi-freeness, with the s-wave weights at the half-integer offset} [T5, Premise 8.1]'s linear assignment the $k = 0$ mode is excluded from this clause and from the converse: at $\omega_0 = 0$ the marginal is $\tfrac12\,\mathrm{id}$ for every $\beta_0$ and constrains nothing).

\end{proof}

\begin{theorem}[Rate rigidity: the temperature is not algebra intrinsic]\label{c:singularity.rate-rigidity-temperature-not-algebra-intrinsic}
Let $\alpha_t = \operatorname{Ad}(e^{itK})$ be the boost flow of the single-mode algebra $\Aa$, let $\theta$ be a $*$-automorphism of $\Aa$, and let $\omega_\beta$ denote the $\beta$-KMS state of $\alpha$.
\begin{enumerate}
\item If $\theta$ commutes with the flow then it fixes every Gibbs state of $K$: $\omega_\beta \circ \theta = \omega_\beta$ for every $\beta$.
\item If $\theta$ reverses the flow, $\theta\circ\alpha_t = \alpha_{-t}\circ\theta$, then $\omega_\beta\circ\theta = \omega_{-\beta}$.
\item Under a constant reparametrisation of the flow, $\alpha_t \mapsto \alpha_{ct}$ with $c > 0$, the $\beta$-KMS state of $\alpha$ is the $(\beta/c)$-KMS state of the reparametrised flow: $\omega_\beta$ is unchanged and its inverse temperature is rescaled.
\item Hence no automorphism intertwining the flow carries $\omega_\beta$ to $\omega_{\beta'}$ for $|\beta'| \neq |\beta|$: $|\beta|$ is an invariant of the pair (state, parametrised flow), and the automorphism group acts on the line of KMS states with orbits $\{\beta, -\beta\}$, so that line carries no automorphism-theoretic structure beyond the orientation flip of (2). A condition that singles out one point of the line is therefore either a stipulation of the state itself or a reference to data outside the pair $(\Aa, \alpha)$; by (3), the value of $|\beta|$ is meaningful only relative to a parametrisation of the flow, and it is that parametrisation, the conversion of the flow parameter into geometric time, that the algebra does not supply.
\end{enumerate}
\end{theorem}
\begin{proof}
Every automorphism of a matrix factor is inner, $\theta = \operatorname{Ad}(U)$. (1) Commutation gives $Ue^{itK}U^{*} = c(t)e^{itK}$ with $|c(t)| = 1$; taking determinants, $c(t)^{2} = 1$, and $c$ is continuous with $c(0) = 1$, so $c \equiv 1$ and $U$ commutes with $K$. As $K = \operatorname{diag}(-\tfrac12, \tfrac12)$ is nondegenerate, $U$ is diagonal and $\operatorname{Ad}(U)$ fixes every diagonal matrix, in particular $e^{-\beta K}/Z$. (2) Flow reversal gives $UKU^{*} = -K$, so $U$ exchanges the spectral projectors of $K$ and $Ue^{-\beta K}U^{*} = e^{+\beta K}$. (3) The reparametrised flow is $\operatorname{Ad}(e^{it(cK)})$, generated by $cK$, and by Theorem \ref{c:singularity.kms-uniqueness-factor} its $\beta'$-KMS state is $e^{-\beta' cK}/Z$, which equals $e^{-\beta K}/Z$ exactly when $\beta' = \beta/c$. (4) By Theorem \ref{c:singularity.kms-uniqueness-factor} the KMS states are exactly the Gibbs states, and an automorphism intertwining the flow either commutes with it or reverses it (it conjugates $K$ to $\pm K$ up to a constant, the spectrum $\{-\tfrac12, \tfrac12\}$ of $K$ leaving no other option), so by (1) and (2) it acts on the parameter through $\{\pm 1\}$ only. That is the whole mathematical content of the clause; the sentence about what a selecting condition must refer to is its reading, and clause (3) is the reason the reading is exact: $\beta$ is the ratio of the state's modular parameter to the flow's parameter, and the second of these is a choice the algebra does not make. The script verifies clauses (1) to (3) on rational data.

\end{proof}

\begin{theorem}[Exact equilibrium over-determines the boundary modes]\label{c:singularity.exact-equilibrium-over-determines-boundary-modes}
Let a state satisfy the geometric modular condition exactly. Then $\varrho \propto e^{-\beta K_{\mathrm{geo}}} = \bigotimes_k e^{-\beta\omega_k K_k}$ is a product across modes, and consequently: every inter-mode correlation vanishes, $C(\varrho) = \sum_k S(\varrho_k) - S(\varrho) = 0$; the descent deficit vanishes, $\rel{\varrho}{\EN\varrho} = 0$ by Theorem \ref{c:singularity.deficit-theorem-cost-descent-exactly}, the state being already a fixed point of $\EN$; and no information whatever is carried by the inter-quanta correlations of Corollary \ref{c:singularity.deficit-carried-quanta-correlations}, the two levels agreeing exactly. The exact geometric modular condition is therefore available only in the stationary idealisation, and the correlations on which Steps 5 and 6b depend are precisely the deviation from it.
\end{theorem}
\begin{proof}
The product form is Theorem \ref{c:singularity.kms-uniqueness-factor} applied to $K_{\mathrm{geo}}$, whose exponential factorises because the $K_k$ commute. A product state has $S(\varrho) = \sum_k S(\varrho_k)$, so $C(\varrho) = 0$. A product of diagonal matrices is diagonal in the joint number basis, hence a fixed point of $\EN$, and the deficit theorem (\ref{c:singularity.deficit-theorem-cost-descent-exactly}) gives $S(\EN\varrho) - S(\varrho) = 0$. The last clause is Corollary \ref{c:singularity.deficit-carried-quanta-correlations}, which identifies the deficit with $C(\varrho) - C(\EN\varrho)$. The script verifies the fixed-point and factorisation clauses in exact arithmetic against an entangled even control failing both.

\end{proof}

\begin{remark}[Two clocks, one calibration]\label{c:note.singularity.two-clocks-one-calibration}
Theorem \ref{c:singularity.rate-rigidity-temperature-not-algebra-intrinsic} is a no-go, and a constructive one: it says exactly what kind of input is needed and how much of it. The inverse temperature is a ratio of two clocks, the modular clock of the state and the geometric clock of the Killing boost, and a ratio of clocks is meaningless until one is calibrated. The chain calibrates exactly once, in Proposition \ref{c:singularity.plane-rate}: an interval of $2\pi$ in rapidity is $2\pi/\kappa = 8\pi M$ of Schwarzschild time. Everything numerical in this step passes through that conversion, as the step already said, and the rigidity theorem now proves that no rearrangement of the algebra can remove the need for it. The residual sign in clause (2) is not new either: it is the orientation of the boost, the labelling of outgoing against ingoing, flagged in Remark \emph{The frame, made explicit} [T5, Rem 6.5] as the $\mathbb{Z}_2$ attached to the weights of $v$ and $v^{*}$. The rigidity theorem identifies that $\mathbb{Z}_2$ as the entire automorphism-theoretic freedom in the temperature.
\end{remark}

\begin{remark}[Why this rescues rather than damages the chain]\label{c:note.singularity.why-rescues-rather-than-damages-chain}
Read carelessly, Theorem \ref{c:singularity.exact-equilibrium-over-determines-boundary-modes} is an obstruction: the state that exact thermality requires is the state with no information in it. Read carefully, it is the reason the order of argument in this subsection matters. Had per-mode thermality been derived from the exact equilibrium condition, as hypothesis (ii) suggested, then thermality and correlation would have been in direct competition and the chain could not have had both. Theorem \ref{c:singularity.ascended-marginals-are-diagonal} shows that per-mode thermality does not rest on equilibrium at all: it rests on parity, which survives an arbitrary non-stationary history. The state class the chain actually needs is
\[
\bigl\{\ \text{even states, per-mode populations at } \beta = 2\pi/\kappa, \ \text{correlations free}\ \bigr\},
\]
in which every single quantum is exactly thermal and the joint state of boundary modes and radiation together is free to carry everything the evaporation writes. Where it is carried is sharpened by the rate model and should be recorded here: under (I) and (L$_E$) the window marginal itself is fixed to be the product member at every stage (Theorem \emph{Eviction: the window is a thermal product, and the radiation entropy is its content} [T5, Thm 11.5](1)), its internal correlation and descent deficit identically zero, and the writing lies across the window/radiation cut and among the quanta of the radiation, not inside the boundary-mode marginal. That this class is where (P) and (R) put the state, and that the dynamics of the boundary modes keeps it there, is Corollary \emph{The chain's state class, and its stability} [T5, Cor 7.32] below; the finite-mode model of the section ``The finite-mode model: the dissolution step by step'' lies in it, and Lemma \emph{Parity protection} [T5, Lem 12.1] is the model's proof that it does. The theorem also settles a question the chain could otherwise have been asked indefinitely: no sharpening of Step 4d will ever supply Step 6b's correlations, because no equilibrium principle can.
\end{remark}

\section{Entropy and information}

\begin{proposition}[Number-diagonal dynamics; the model's scrambling derived]\label{c:singularity.number-diagonal-dynamics-models-scrambling-derived}
Under (L$_E$) every generator available to the internal dynamics of the boundary modes lies in the abelian algebra generated by $\{\sigma_k\}$, hence is diagonal in the joint number basis and its exponential is a diagonal unitary. Consequently each mode's occupation probabilities are constants of the motion of the boundary modes, so populations set at injection by the local geometry do not drift; and the diagonal scrambling of the finite-mode model, the controlled-$Z$ circuit of the section ``The model'', is not a modelling convenience but the general form of what (L$_E$) permits.
\end{proposition}
\begin{proof}
The $\sigma_k = 2N_k - 1$ commute and are simultaneously diagonal in the joint number basis; a polynomial in commuting diagonal operators is diagonal and its exponential is a diagonal unitary, which multiplies each joint-number amplitude by a phase and so leaves every joint-number probability, hence every marginal population, unchanged. With Theorem \ref{c:singularity.ascended-marginals-are-diagonal} fixing the off-diagonal entries at zero, each $\varrho_k$ is constant in time. The script verifies the invariance for a controlled-$Z$ circuit on the model's mode space and exhibits an orthogonal, parity-preserving hopping rotation that does change the marginal populations, which is exactly the generator (L$_E$, \ref{c:singularity.l-boundary-mode-locality}) excludes.

\end{proof}

\section{Results used from elsewhere}
The proofs above use the following results of the series, stated briefly; each is proved in the paper named.

\begin{fact}[The curvature of the gauge field: exact shape decoupling]\label{c:singularity.curvature-gauge-field-exact-shape-decoupling}
The curvature of the black-hole gauge field is exactly $\mathcal{R} = -(M/r^3)\,S$ with $S = 1 - 3\Pi$ a fixed shape operator of spectrum $(-2,-2,1,1,1,1)$.
\emph{Source:} T5, Thm 6.1 (\href{https://doi.org/10.5281/zenodo.22721790}{doi:10.5281/zenodo.22721790}).
\end{fact}

\begin{fact}[The deficit is carried by the quanta correlations]\label{c:singularity.deficit-carried-quanta-correlations}
When every single-mode marginal is number-diagonal, the descent deficit $\rel{\rho}{\EN\rho}$ equals the drop in total correlation $C(\rho) - C(\EN\rho)$.
\emph{Source:} T5, Cor 9.4 (\href{https://doi.org/10.5281/zenodo.22721790}{doi:10.5281/zenodo.22721790}).
\end{fact}

\begin{fact}[The deficit theorem: the cost of descent, exactly]\label{c:singularity.deficit-theorem-cost-descent-exactly}
For any pinching $E$ and density matrix $\rho$, $S(E(\rho)) - S(\rho) = \rel{\rho}{E(\rho)} \geq 0$, with equality exactly when $\rho = E(\rho)$.
\emph{Source:} T5, Thm 9.3 (\href{https://doi.org/10.5281/zenodo.22721790}{doi:10.5281/zenodo.22721790}).
\end{fact}

\begin{fact}[The descent, and what it erases]\label{c:singularity.descent-what-it-erases}
The descent $E(x) = PxP + (1-P)x(1-P)$, with $P = v^* v$, keeps a mode's diagonal occupation data and erases its coherences.
\emph{Source:} T5, Prop 5.5 (\href{https://doi.org/10.5281/zenodo.22721790}{doi:10.5281/zenodo.22721790}).
\end{fact}

\begin{fact}[The Euclidean rotor, and where it lies]\label{c:singularity.euclidean-rotor-where-it-lies}
The imaginary-rapidity continuation $R_{i\theta} = \cos(\theta/2) + i\sigma\sin(\theta/2)$ of the radial boost rotor has period $4\pi$, with conjugation period $2\pi$, and is realised by no real rotor of $\Cl(1,3)$, in particular not by the curvature-selected pair, whose dual member $I\sigma_r = e_2e_3$ rotates along the bifurcation sphere.
\emph{Source:} T5, Prop 7.21 (\href{https://doi.org/10.5281/zenodo.22721790}{doi:10.5281/zenodo.22721790}).
\end{fact}

\begin{fact}[Modular theory of a matrix factor, proved outright]\label{c:singularity.modular-theory-matrix-factor-proved-outright}
For a full matrix algebra with a cyclic separating vector, the modular operator, modular conjugation and modular flow are computed explicitly, and the state is KMS for that flow.
\emph{Source:} T5, Lem 7.1 (\href{https://doi.org/10.5281/zenodo.22721790}{doi:10.5281/zenodo.22721790}).
\end{fact}

\begin{fact}[The plane and the rate]\label{c:singularity.plane-rate}
The curvature selects the boost plane of the completed pair invariantly, and at the horizon its rate is the surface gravity $\kappa = 1/4M$, so $2\pi$ of rapidity corresponds to $8\pi M$ of Schwarzschild time.
\emph{Source:} T5, Prop 6.2 (\href{https://doi.org/10.5281/zenodo.22721790}{doi:10.5281/zenodo.22721790}).
\end{fact}

\begin{thebibliography}{99}
\bibitem{ButlerT5} L.~Butler,
\emph{Division by Zero at the Gravitational Singularity: The Transport Principle Applied to the Black Hole, and the Information Paradox as a Category Error},
Transport Papers T5 (2026), doi:\href{https://doi.org/10.5281/zenodo.22721790}{10.5281/zenodo.22721790}.
\bibitem{Hawking1975} S.~W.~Hawking,
Particle creation by black holes,
\emph{Commun.\ Math.\ Phys.} \textbf{43} (1975), 199--220.
\bibitem{HartleHawking1976} J.~B.~Hartle and S.~W.~Hawking,
Path-integral derivation of black-hole radiance,
\emph{Phys.\ Rev.\ D} \textbf{13} (1976), 2188--2203.
\bibitem{Unruh1976} W.~G.~Unruh,
Notes on black-hole evaporation,
\emph{Phys.\ Rev.\ D} \textbf{14} (1976), 870--892.
\bibitem{Boulware1975} D.~G.~Boulware,
Quantum field theory in Schwarzschild and Rindler spaces,
\emph{Phys.\ Rev.\ D} \textbf{11} (1975), 1404--1423.
\bibitem{KayWald1991} B.~S.~Kay and R.~M.~Wald,
Theorems on the uniqueness and thermal properties of stationary, nonsingular, quasifree states on spacetimes with a bifurcate Killing horizon,
\emph{Phys.\ Rep.} \textbf{207} (1991), 49--136.
\bibitem{FredenhagenHaag1990} K.~Fredenhagen and R.~Haag,
On the derivation of Hawking radiation associated with the formation of a black hole,
\emph{Commun.\ Math.\ Phys.} \textbf{127} (1990), 273--284.
\bibitem{ButlerT4} L.~Butler,
\emph{Clifford Sheaves over Degenerate Metrics: Nilpotent Ideals, One-Sided Completion, and the Exceptional Algebra},
Transport Papers T4 (2026), doi:\href{https://doi.org/10.5281/zenodo.22721669}{10.5281/zenodo.22721669}.
\bibitem{ButlerT3} L.~Butler,
\emph{The Tower and the CAR Limit: Iterated Transport, the Frozen Clock, and Second Quantisation},
Transport Papers T3 (2026), doi:\href{https://doi.org/10.5281/zenodo.22721556}{10.5281/zenodo.22721556}.
\bibitem{BR2} O.~Bratteli and D.~W.~Robinson,
\emph{Operator Algebras and Quantum Statistical Mechanics 2: Equilibrium States. Models in Quantum Statistical Mechanics},
2nd ed., Springer (1997).
\bibitem{HHW1967} R.~Haag, N.~M.~Hugenholtz and M.~Winnink,
On the equilibrium states in quantum statistical mechanics,
\emph{Commun.\ Math.\ Phys.} \textbf{5} (1967), 215--236.
\bibitem{Connes1973} A.~Connes,
Une classification des facteurs de type III,
\emph{Ann.\ Sci.\ \'Ec.\ Norm.\ Sup\'er.} (4) \textbf{6} (1973), 133--252.
\bibitem{GibbonsHawking1977} G.~W.~Gibbons and S.~W.~Hawking,
Action integrals and partition functions in quantum gravity,
\emph{Phys.\ Rev.\ D} \textbf{15} (1977), 2752--2756.
\bibitem{Israel1976} W.~Israel,
Thermo-field dynamics of black holes,
\emph{Phys.\ Lett.\ A} \textbf{57} (1976), 107--110.
\bibitem{BW1976} J.~J.~Bisognano and E.~H.~Wichmann,
On the duality condition for quantum fields,
\emph{J.\ Math.\ Phys.} \textbf{17} (1976), 303--321.
\bibitem{Sewell1982} G.~L.~Sewell,
Quantum fields on manifolds: PCT and gravitationally induced thermal states,
\emph{Ann.\ Phys.} \textbf{141} (1982), 201--224.
\bibitem{ButlerVerifySS} L.~Butler,
\emph{verify\_state\_selection.py: exact-arithmetic verification of the computational identities of ``Selecting the State''},
Python script, deposited with this note alongside the shared library GA.py (2026).
\bibitem{WWW1952} G.~C.~Wick, A.~S.~Wightman and E.~P.~Wigner,
The intrinsic parity of elementary particles,
\emph{Phys.\ Rev.} \textbf{88} (1952), 101--105.
\end{thebibliography}
\end{document}